

\documentclass[12pt]{article}

\usepackage{sbc-template}

\usepackage{graphicx,url}
\usepackage[T1]{fontenc}
\usepackage[utf8]{inputenc}
\usepackage[brazil]{babel}
\usepackage{array}
\usepackage{booktabs}
\usepackage{enumitem}

\graphicspath{{figuras/}}


\newenvironment{referencias}{%
  \begin{list}{}{%
    \setlength{\leftmargin}{1.25cm}%
    \setlength{\itemindent}{-1.25cm}%
    \setlength{\labelwidth}{0pt}%
    \setlength{\labelsep}{0pt}%
    \setlength{\itemsep}{6pt}%
    \setlength{\parsep}{0pt}%
    \setlength{\topsep}{6pt}}}%
  {\end{list}}


\newcommand{\fonte}[1]{%
  \vspace{2pt}\par\noindent{\footnotesize Fonte: #1}}

\sloppy

\title{Eärendil: Um Sistema Baseado em \emph{Machine Learning} para
Detecção de Falhas em Pavimentos Urbanos por Meio de Drones}

\author{Fabrício José\inst{1}, Marco Abbud\inst{1}}

\address{Universidade Federal de Ciência e Inovação do Rio de Janeiro\\
  Rio de Janeiro, RJ, Brasil
  \email{fabricio.macedo@aluno.cefet-rj.br, marco.abbud@cefet-rj.br}
  % TODO: ajustar o e-mail do orientador
}

\begin{document}

\maketitle


\begin{abstract}
  This work proposes Earendil, an integrated, low-cost and open-source
  system for automatic pothole detection on urban pavements from
  drone-captured images. The system comprises four components: an
  external image source (drone), a backend responsible for ingestion,
  normalization and persistence, a YOLO26-based inference service, and a
  web interface that displays georeferenced records on a map. The
  detector is trained on a unified dataset built from the public
  UAV-PDD2023 and RDD2022 benchmarks, restricted to aerial capture, and
  evaluated with standard object detection metrics.
\end{abstract}

\begin{resumo}
  Este trabalho propõe o Earendil, um sistema integrado, de baixo custo e
  baseado em componentes de código aberto para a detecção automática de
  buracos em pavimentos urbanos a partir de imagens capturadas por drone.
  O sistema é composto por quatro componentes: a fonte externa de imagens
  (drone), um \emph{backend} responsável pela ingestão, normalização e
  persistência dos registros, um serviço de inferência baseado na
  arquitetura YOLO26 e uma interface web que apresenta os registros
  georreferenciados em um mapa. O detector é treinado sobre um conjunto
  unificado construído a partir das bases públicas UAV-PDD2023 e
  RDD2022, restrito ao domínio de captura aérea, e avaliado por métricas
  consagradas de detecção de objetos.
\end{resumo}

% =====================================================================
\section{Introdução}
\label{sec:introducao}

A malha rodoviária é o principal meio de circulação de pessoas e de
cargas no Brasil, e sua conservação tem efeito direto sobre a segurança
dos usuários e sobre o custo do transporte. A Pesquisa CNT de Rodovias
2025 avaliou 114.197 km de vias pavimentadas e classificou 62,1\% da
extensão como regular, ruim ou péssima em seu estado geral; segundo o
mesmo levantamento, a qualidade do pavimento eleva em média 31,2\% os
custos operacionais do transporte rodoviário no país (CNT, 2025). O
quadro é ainda mais desfavorável nas rodovias sob gestão pública, em que
64,4\% da extensão apresenta algum problema de pavimento.

Nas vias urbanas, sob responsabilidade de municípios que frequentemente
dispõem de orçamento e corpo técnico reduzidos, o problema é agravado
pela ausência de levantamentos periódicos e padronizados equivalentes aos
realizados na malha federal. Entre as patologias de pavimentos flexíveis,
o buraco, também denominado panela na terminologia normativa (DNIT,
2003), é a que produz efeito mais imediato sobre o usuário: representa
estágio avançado de degradação, originado em geral pela evolução de
trincas de fadiga e pela infiltração de água, e está associado a danos em
veículos, a desvios bruscos de trajetória e a acidentes, sobretudo com
motociclistas e ciclistas.

A avaliação da condição de superfície é hoje conduzida majoritariamente
por levantamento visual, procedimento previsto na norma DNIT
008/2003-PRO, no qual técnicos treinados percorrem os trechos registrando
tipo, severidade e extensão dos defeitos. Embora de baixo custo unitário,
o método é lento, depende do julgamento do avaliador, expõe o inspetor ao
tráfego e muitas vezes exige interdição parcial da via. Na prática
municipal, esse processo convive com um canal reativo e enviesado: a
reclamação do cidadão, que registra apenas os defeitos percebidos por
quem decide reclamar e raramente traz evidência visual ou localização
confiável.

O uso de veículos aéreos não tripulados (VANTs), popularmente chamados de
drones, associado a técnicas de visão computacional, apresenta-se como
alternativa promissora. O drone cobre grandes áreas em pouco tempo,
alcança trechos de difícil acesso, dispensa a interrupção do tráfego e
não expõe o operador ao risco de atropelamento; além disso, registra a
posição de captura por meio de GPS embarcado. Sobre essas imagens,
detectores de objetos em tempo real da família YOLO (\emph{You Only Look
Once}), que processam a imagem inteira em uma única passagem pela rede
(Redmon et al., 2016), permitem localizar automaticamente os defeitos com
custo computacional compatível com \emph{hardware} modesto.

Há, no entanto, uma lacuna recorrente na literatura nacional sobre o
tema: a maior parte dos trabalhos encerra-se no modelo. Treina-se um
detector, relatam-se métricas de precisão e revocação em um conjunto de
teste e conclui-se pela viabilidade da abordagem, sem que se construa o
caminho entre a imagem capturada em campo e a informação efetivamente
utilizável por quem decide onde aplicar o recurso de manutenção. Faltam a
ingestão automatizada das imagens, a persistência dos registros, o
encadeamento com o modelo de inferência e a apresentação dos resultados
em uma interface consultável.

Este trabalho ataca precisamente essa lacuna. A questão que orienta a
pesquisa é: \emph{é possível construir um sistema integrado, de baixo
custo e baseado em componentes de código aberto, capaz de receber imagens
de pavimento capturadas por drone, classificar automaticamente a presença
de buracos e disponibilizar os registros localizados em uma interface web
de consulta, com desempenho de detecção compatível com o relatado na
literatura?} A proposta apresentada a seguir, denominada Earendil,
responde a essa questão por meio de um sistema de quatro componentes e de
um detector treinado sobre dois conjuntos de dados públicos de
referência.

% =====================================================================
\section{Solução Proposta}
\label{sec:solucao}

\subsection{Objetivo geral}
\label{sec:objetivo-geral}

Projetar, implementar e avaliar o Earendil, um sistema baseado em
aprendizado de máquina para detecção automática de buracos em pavimentos
urbanos a partir de imagens capturadas por drone, contemplando o fluxo
completo de captura, processamento, persistência, inferência e
visualização georreferenciada dos registros.

\subsection{Objetivos específicos}
\label{sec:objetivos-especificos}

\begin{itemize}
  \item Construir um conjunto de dados de treinamento unificado a partir
    dos \emph{datasets} públicos UAV-PDD2023 e RDD2022, harmonizando
    taxonomias de rótulos e formatos de anotação.
  \item Treinar e avaliar modelos de detecção baseados na arquitetura
    YOLO26, comparando variantes de tamanho e estabelecendo uma linha de
    base com versões anteriores da família YOLO.
  \item Avaliar a capacidade de generalização do modelo entre domínios,
    treinando em um \emph{dataset} e testando no outro, de modo a estimar
    o desempenho esperado em imagens não vistas.
  \item Implementar um serviço de \emph{backend} responsável por receber
    as imagens enviadas pela fonte externa, normalizá-las, registrá-las
    em banco de dados e acionar o modelo de inferência.
  \item Implementar uma interface web que apresente os registros de
    detecção com a respectiva imagem anotada, nível de confiança, data e
    posição aproximada obtida do GPS do drone.
  \item Avaliar o sistema resultante por meio de testes unitários, de
    integração e de usabilidade, e o modelo por meio de métricas
    consagradas de detecção de objetos.
\end{itemize}

\subsection{Arquitetura proposta}
\label{sec:arquitetura}

O Earendil é organizado em quatro componentes com responsabilidades bem
delimitadas, comunicando-se por interfaces explícitas. Essa separação
permite que o modelo seja substituído ou re-treinado sem alteração dos
demais módulos e que a fonte de imagens seja trocada sem impacto no
restante do sistema.

\textbf{Fonte externa (drone).} Responsável por registrar as imagens do
pavimento e enviá-las ao \emph{backend}. A captura ocorre em voo a baixa
altitude sobre a via, com a câmera apontada para a superfície; cada
imagem é acompanhada das coordenadas de latitude e longitude fornecidas
pelo GPS do equipamento e do instante da captura. O envio é feito por
requisição HTTP ao \emph{endpoint} de ingestão do \emph{backend}. Como o
objetivo do trabalho é o sistema, e não a automação do voo, esse
componente será exercitado tanto por um agente de envio em lote que
percorre um diretório de imagens reais de drone e as submete ao
\emph{backend} com os metadados correspondentes quanto, quando houver
disponibilidade de equipamento, por captura real em campo.

\textbf{Backend.} Núcleo de orquestração do sistema. Recebe a imagem e
seus metadados, normaliza o arquivo, atribuindo-lhe um identificador
único e um nome padronizado que elimina colisões e preserva a
rastreabilidade, logo em seguida armazena o binário no repositório de
arquivos e cria o registro correspondente no banco de dados com estado
inicial \emph{pendente}. Em seguida, submete a imagem ao modelo de
aprendizado de máquina e persiste o resultado: classe atribuída, nível de
confiança, caixas delimitadoras e estado final do registro
(\emph{confirmado} ou \emph{descartado}). Por fim, expõe os registros por
meio de uma API de consulta, consumida pelo \emph{frontend}.

\textbf{Modelo de aprendizado de máquina.} Responsável por classificar a
imagem do pavimento, indicando se há buraco e, em caso afirmativo, em que
região da imagem. Será implementado com a arquitetura YOLO26, treinada
por transferência de aprendizado a partir de pesos pré-treinados e
ajustada sobre o conjunto unificado descrito na
Seção~\ref{sec:dados}. O modelo é exposto ao \emph{backend} como um
serviço de inferência, recebendo a imagem e devolvendo a lista de
detecções em formato estruturado.

\textbf{Frontend.} Interface web responsável por dar visibilidade aos
registros produzidos. Apresenta um mapa com marcadores nas posições
informadas pelo GPS do drone e uma listagem tabular filtrável por data,
classe e nível de confiança. Ao selecionar um registro, o usuário
visualiza a imagem com as caixas delimitadoras sobrepostas e os metadados
associados. A localização é tratada de forma aproximada: adota-se
diretamente a coordenada registrada pelo drone, sem correção geométrica
da projeção nem ortorretificação, o que é suficiente para o propósito de
indicar em que trecho da via o defeito se encontra.

\subsection{Fluxo de execução}
\label{sec:fluxo}

\begin{enumerate}
  \item O drone sobrevoa a via e captura imagens da superfície do
    pavimento, associando a cada imagem as coordenadas do GPS e o
    instante da captura.
  \item A fonte externa submete a imagem e os metadados ao
    \emph{endpoint} de ingestão do \emph{backend}.
  \item O \emph{backend} valida a requisição, renomeia e armazena a
    imagem e cria o registro no banco de dados com estado
    \emph{pendente}.
  \item O \emph{backend} envia a imagem ao modelo, que retorna as
    detecções encontradas com as respectivas confianças.
  \item O \emph{backend} atualiza o registro com o resultado da
    inferência, marcando-o como \emph{confirmado} quando há detecção
    acima do limiar de confiança e como \emph{descartado} caso contrário.
  \item O \emph{frontend} consulta a API e apresenta os registros no mapa
    e na listagem, permitindo a inspeção individual de cada ocorrência.
\end{enumerate}

\subsection{Contribuições esperadas}
\label{sec:contribuicoes}

A principal contribuição é de natureza sistêmica: entregar, de ponta a
ponta, o encadeamento entre captura aérea, detecção automática e
visualização, em vez de restringir o trabalho ao treinamento de um
detector. Espera-se, adicionalmente, produzir três resultados de
interesse: um conjunto de dados unificado a partir de duas bases públicas
de referência, com taxonomia harmonizada e documentada; uma avaliação da
arquitetura YOLO26, cuja inferência nativamente livre de supressão de
não-máximos simplifica a integração com o \emph{backend}, aplicada a
imagens aéreas de pavimento, tarefa para a qual ainda há poucos relatos
publicados; e um protótipo funcional, licenciado em código aberto,
reaproveitável por órgãos municipais e por trabalhos subsequentes.

% =====================================================================
\section{Trabalhos Relacionados}
\label{sec:trabalhos-relacionados}

O levantamento a seguir prioriza trabalhos que abordam a detecção
automática de defeitos em pavimentos por visão computacional, com atenção
especial aos que utilizam imagens aéreas e aos que foram publicados no
contexto brasileiro. Ao final, a Tabela~\ref{tab:comparacao} sintetiza a
comparação e a Seção~\ref{sec:delimitacoes} delimita o que a proposta
deliberadamente não pretende contemplar.

\subsection{Análise comparativa}
\label{sec:analise-comparativa}

O trabalho mais próximo da proposta é o de Alzamzami et al. (2024), que
apresentam o PDS-UAV, sistema de detecção e gestão de buracos baseado em
imagens de drone. Os autores conduziram entrevista com servidor da
prefeitura de Jeddah e questionário com 101 usuários das vias para
levantar requisitos, modelaram o sistema em UML e implementaram duas
aplicações web: uma para os servidores da manutenção viária, que
cadastram e atualizam o estado dos buracos, e outra para os usuários da
via, que consultam as ocorrências ao longo de uma rota. O detector foi
treinado com 1.362 imagens públicas capturadas por drone na Espanha e
testado em 111 imagens de Jeddah, em um experimento de generalização
entre cidades; entre YOLOv4-tiny, YOLOv5 e YOLOv8, a última obteve o
melhor resultado, com precisão de 98\%, revocação de 92\% e F1 de 95\%. A
proximidade com o Earendil é evidente na arquitetura (drone, modelo,
banco e aplicação web), e o PDS-UAV é adotado aqui como principal
referência de projeto.

As diferenças, contudo, são relevantes. No PDS-UAV, o envio das imagens é
uma ação manual do servidor, que preenche um formulário com coordenadas e
distrito antes de acionar a predição; no Earendil, a ingestão é
automatizada por uma interface programática, e a coordenada vem do
próprio GPS do equipamento. Enquanto o PDS-UAV utiliza YOLOv8, esta
proposta avalia o YOLO26, arquitetura cuja saída direta, sem etapa de
pós-processamento por supressão de não-máximos, reduz a latência e
simplifica a integração em serviço (Ultralytics, 2026). Por fim, o
treinamento do PDS-UAV recorre a um único \emph{dataset} de um só país,
ao passo que o Earendil combina uma base aérea multiclasse com uma base
multinacional, buscando maior robustez a variações de iluminação, textura
e ângulo de captura.

Monte, Goldschmidt e Santos (2025) propõem um modelo híbrido baseado em
YOLOv8 para identificação e classificação de patologias em pavimentos
flexíveis e rígidos, área em que apontam lacuna após revisão sistemática
de 209 estudos publicados entre 2019 e 2024. Coletaram cerca de 4.000
imagens com câmeras GoPro equipadas com GPS em diferentes regiões do
Brasil e sob condições climáticas variadas, anotando 1.000 delas
manualmente, e relatam mAP@50 próximo de 90\% e precisão de até 96\%, com
validação adicional em 65 segmentos de um quilômetro. É o trabalho
brasileiro com melhores métricas entre os analisados e reforça duas
escolhas desta proposta: a importância da diversidade de condições de
captura e o valor da anotação cuidadosa. A captura, no entanto, é
veicular e em nível de solo, e o trabalho concentra-se no modelo, sem
descrever a entrega de um sistema de consulta; o \emph{dataset} é
proprietário, o que dificulta a reprodução.

Carvalho, Guedes e Figueiredo (2022) tratam a inspeção automática de
danos em pavimentação como problema de detecção em visão computacional,
no contexto de cidades inteligentes. Avaliam quatro configurações de
redes YOLO sobre dados de três países e obtêm mAP de 0,53, com ganhos
expressivos no tempo de predição em relação à literatura da época, ainda
que com menor número de parâmetros. O estudo de caso conduzido com
imagens brasileiras é particularmente instrutivo para esta proposta, pois
evidencia a queda de desempenho quando o modelo é aplicado a um domínio
distinto do de treinamento, justamente o risco que o teste de
generalização entre \emph{datasets}, previsto na
Seção~\ref{sec:treinamento}, pretende quantificar.

Batista et al. (2025) avaliam o YOLOv8n sobre um conjunto real de imagens
de vias urbanas do Ceará, anotadas com uma única classe, ``defeito'', na
plataforma Roboflow. Obtêm precisão de 74,60\%, revocação de 44,18\% e
mAP@50 de 0,521, resultado que consideram satisfatório para ocorrências
visuais simples. O trabalho é útil como referência de desempenho em
contexto urbano brasileiro e ilustra o custo de uma classe única
excessivamente ampla: a revocação baixa indica que boa parte dos defeitos
presentes não é recuperada. O Earendil concentra-se em uma classe
específica e visualmente bem caracterizada, o buraco, na expectativa de
obter revocação mais alta.

Wanderley e Silva (2025) investigam a identificação de patologias em
pavimentos a partir de imagens de drone com modelos YOLO treinados no
Roboflow e executados em Google Colab. Desenvolveram cinco modelos, dos
quais quatro foram avaliados: o mais consistente, com 909 imagens de
treino e curadoria manual intensiva, alcançou mAP@50 de 55,8\%, precisão
de 60,6\% e revocação de 56,4\%; já um modelo treinado de forma
automatizada, sem verificação manual das anotações e com base 2,5 vezes
maior, caiu para mAP@50 de 29,5\%. A conclusão é que a qualidade das
anotações é mais determinante que o volume bruto de dados, e portanto é a
principal lição incorporada por esta proposta e justifica a decisão de
partir de \emph{datasets} de referência revisados por pares, com
anotações validadas, em vez de compor uma base própria por coleta
oportunista.

Kassem, Matrakas e Bussador (2024) propõem um protótipo baseado em YOLO
para detecção de rachaduras em concreto asfáltico, mas o trabalho
permanece em estágio de proposta: a construção do banco de imagens, o
treinamento e a avaliação são apontados como etapas futuras. O contraste
é deliberado nesta seção, pois delimita um risco concreto a ser evitado,
o de encerrar o trabalho no desenho da solução. O cronograma da
Seção~\ref{sec:cronograma} reserva, por isso, blocos explícitos de
treinamento, integração e validação.

Guan et al. (2025) apresentam um sistema embarcado de detecção de buracos
que combina aprendizado profundo e visão binocular. Constroem um
\emph{dataset} de 6.848 imagens com câmera binocular veicular, propõem o
modelo GASB-ShuffleNetv2 para reconhecimento do estado do pavimento com
acurácia de 98,61\% e o PF-CenterNet para localização das regiões
afetadas, além de um modelo de estimativa de distância por disparidade
estéreo, executado em um Jetson Xavier NX com aceleração por TensorRT e
tempo total de 111 ms por imagem. O objetivo é distinto do desta
proposta: trata-se de percepção em tempo real para alerta ao condutor, e
não de apoio à gestão do patrimônio viário. O trabalho é citado por
demonstrar que a inferência pode ser executada em \emph{hardware}
embarcado de baixo consumo, o que sustenta a escolha de uma arquitetura
orientada à borda como o YOLO26 e abre caminho para uma evolução futura
do Earendil com inferência a bordo do próprio drone.

\begin{table}[ht]
\centering
\caption{Comparação entre os trabalhos relacionados e a proposta}
\label{tab:comparacao}
\footnotesize
\setlength{\tabcolsep}{3pt}
\renewcommand{\arraystretch}{1.15}
\begin{tabular}{p{2.4cm} p{2.0cm} p{1.9cm} p{3.2cm} p{3.6cm}}
\hline
\textbf{Trabalho} & \textbf{Captura} & \textbf{Algoritmo} &
\textbf{Dados} & \textbf{Entrega / resultado} \\
\hline
Alzamzami et al. (2024): PDS-UAV &
Drone (DJI) &
YOLOv8 &
Público (Espanha, 1.362 img.) + 111 img. de Jeddah &
Sistema web + modelo; P = 98\%, R = 92\%, F1 = 95\% \\
\hline
Monte et al. (2025) &
Câmera GoPro com GPS em veículo &
YOLOv8 híbrido &
Próprio, $\sim$4.000 img. (1.000 anotadas) &
Modelo; mAP@50 $\approx$ 90\%, P até 96\% \\
\hline
Carvalho et al. (2022) &
Câmera veicular &
YOLO (4 config.) &
RDD (3 países) + estudo de caso Brasil &
Modelo; mAP $\approx$ 0,53 \\
\hline
Batista et al. (2025) &
Câmera em nível de solo &
YOLOv8n &
Próprio (vias do Ceará), classe única &
Modelo; P = 74,6\%, R = 44,2\%, mAP@50 = 0,521 \\
\hline
Wanderley e Silva (2025) &
Drone / imagens diversas &
YOLO (Roboflow) &
Roboflow Universe, $\sim$1.048 img. &
Modelo; mAP@50 = 55,8\%, P = 60,6\%, R = 56,4\% \\
\hline
Kassem et al. (2024) &
Imagens digitais de asfalto &
YOLO (proposto) &
Não construído &
Proposta metodológica, sem implementação \\
\hline
Guan et al. (2025) &
Câmera binocular embarcada &
GASB-Shuffle\-Netv2 + PF-CenterNet &
Próprio, 6.848 img. &
Sistema embarcado com estimativa de distância; acurácia 98,6\% \\
\hline
\textbf{Earendil (proposta)} &
Drone (imagem + GPS) &
YOLO26 &
UAV-PDD2023 + RDD2022 (públicos) &
Sistema completo (ingestão, banco, inferência, mapa web) + avaliação \\
\hline
\end{tabular}
\fonte{elaborado pelo autor, 2026.}
\end{table}

Em síntese, os trabalhos revisados consolidam a viabilidade técnica da
detecção automática de buracos por redes da família YOLO, com métricas
que variam de mAP@50 em torno de 0,52 em cenários urbanos de classe única
até precisão superior a 95\% em bases mais controladas. A lacuna que o
Earendil ocupa não está, portanto, na detecção em si, mas na integração:
entre os trabalhos analisados, apenas o PDS-UAV entrega um sistema
operacional baseado em imagens de drone, e mesmo ele depende de submissão
manual das imagens e não explora bases de dados aéreas multiclasse nem
arquiteturas de detecção ponta a ponta.

\subsection{Delimitações do escopo}
\label{sec:delimitacoes}

Para que o trabalho seja concluído no prazo previsto e com profundidade
adequada, os seguintes aspectos ficam explicitamente fora do escopo:

\begin{itemize}
  \item \textbf{Georreferenciamento preciso e ortorretificação.} A
    posição de cada registro será a coordenada bruta informada pelo GPS
    do drone, sem correção da projeção da imagem, sem geração de
    ortomosaico e sem uso de pontos de controle no solo. O resultado é
    uma localização aproximada, adequada para indicar o trecho, não para
    medir a posição exata do defeito na pista.
  \item \textbf{Quantificação de severidade e de geometria.} O sistema
    indicará a presença e a posição do buraco na imagem, mas não estimará
    profundidade, área real em metros quadrados nem grau de severidade
    segundo a norma DNIT 005/2003-TER, o que exigiria reconstrução
    tridimensional ou escala conhecida.
  \item \textbf{Detecção do conjunto completo de patologias.} O trabalho
    concentra-se na classe buraco. Trincas, afundamentos, exsudação e
    desagregação superficial poderão ser aproveitadas como classes
    auxiliares no treinamento, mas não serão objeto de avaliação nem de
    exibição no \emph{frontend}.
  \item \textbf{Inferência a bordo e planejamento autônomo de voo.} A
    inferência ocorre no servidor, e não no drone; o planejamento de
    rota, o controle do voo e a conformidade com a regulamentação de
    operação de VANTs não são tratados.
  \item \textbf{Integração com sistemas de gestão municipal e fluxo de
    ordens de serviço.} O \emph{frontend} dará visibilidade aos
    registros, sem controle de estado de reparo, atribuição de equipes ou
    integração com sistemas legados de prefeituras.
  \item \textbf{Operação em escala de produção.} Não se pretende
    dimensionar o sistema para uso municipal contínuo, com alta
    disponibilidade e processamento de grandes volumes; o resultado é um
    protótipo funcional validado em escala experimental.
\end{itemize}

% =====================================================================
\section{Metodologia}
\label{sec:metodologia}

\subsection{Processo de desenvolvimento}
\label{sec:processo}

O projeto será conduzido por um processo iterativo e incremental,
organizado em ciclos de duas semanas. Cada ciclo se encerra com um
incremento executável e verificável, o que permite antecipar riscos
técnicos, em especial os associados ao desempenho do modelo, e replanejar
sem comprometer o cronograma global. A gestão das atividades será feita
com auxílio da plataforma de versionamento de código GitHub, e o controle
de versão em repositório Git único contendo os módulos de \emph{backend},
\emph{frontend} e treinamento, que poderá ser encontrado no endereço:
\url{https://github.com/JfOliveira2003/Earendil}

A trilha de aprendizado de máquina seguirá, em paralelo, as etapas do
CRISP-DM: entendimento do problema e dos dados, preparação, modelagem,
avaliação e implantação. A separação das duas trilhas é intencional, pois
seus ciclos têm durações distintas (um experimento de treinamento pode
levar horas, enquanto um incremento de interface se conclui em minutos) e
sua sincronização ocorre nos pontos de integração previstos no
cronograma.

\subsection{Levantamento de requisitos e modelagem}
\label{sec:requisitos}

O levantamento de requisitos combinará três fontes. A primeira é a
análise documental das normas técnicas pertinentes, em especial a DNIT
005/2003-TER, que fixa a terminologia dos defeitos, e a DNIT
008/2003-PRO, que descreve o procedimento de levantamento visual, que são
insumos necessários para que a nomenclatura adotada no sistema seja
compatível com a prática de engenharia. A segunda é a análise de sistemas
similares, tomando o PDS-UAV como referência de funcionalidades e
organização. A terceira, condicionada à disponibilidade, é uma entrevista
semiestruturada com servidor de secretaria municipal de obras ou de
infraestrutura, no formato adotado por Alzamzami et al. (2024), destinada
a caracterizar o processo atual de identificação e registro de buracos e
a validar a utilidade das informações que o \emph{frontend} pretende
apresentar.

A modelagem utilizará notação UML, com os seguintes artefatos: diagrama
de casos de uso, delimitando as funcionalidades e os atores (operador de
captura e consultor dos registros); diagrama de componentes, explicitando
as fronteiras entre os quatro módulos e as interfaces que os conectam;
diagrama de sequência do fluxo principal, da submissão da imagem à
exibição do registro no mapa; e modelo entidade-relacionamento do banco
de dados. Os casos de uso serão descritos em formato textual, com
pré-condições, fluxo principal e fluxos alternativos, e servirão de base
para a definição dos casos de teste.

\subsection{Obtenção e preparação dos dados}
\label{sec:dados}

O treinamento utilizará dois conjuntos de dados públicos, escolhidos por
serem revisados por pares, terem anotações validadas e se complementarem
em perspectiva de captura. O UAV-PDD2023 (Yan e Zhang, 2023) reúne 2.440
imagens capturadas por drone na China, com mais de 11.150 instâncias
anotadas em seis classes de patologia: trincas longitudinais,
transversais, oblíquas e interligadas, remendos e buracos, obtidas com
câmera DJI Zenmuse H20 em voo a cerca de 30 m de altitude, sob duas
condições climáticas e em rodovias de diferentes padrões construtivos,
com anotações no formato Pascal VOC. É a base que mais se aproxima da
geometria de captura prevista para o Earendil. O RDD2022 (Arya et al.,
2024) reúne 47.420 imagens de seis países: Japão, Índia, República
Tcheca, Noruega, Estados Unidos e China. Com mais de 55.000 instâncias de
dano anotadas em quatro classes, entre elas buraco; o conjunto inclui um
subconjunto capturado por drone na China, além de imagens veiculares. Sua
contribuição é a diversidade geográfica e de condições de captura, fator
que Carvalho et al. (2022) associam diretamente à capacidade de
generalização do modelo.

A preparação envolverá quatro etapas. Primeiro, a harmonização de
taxonomias: as classes dos dois conjuntos serão mapeadas para um
vocabulário comum, com a classe buraco como alvo principal. Segundo, a
conversão de formatos: as anotações Pascal VOC do UAV-PDD2023 e as do
RDD2022 serão convertidas para o formato de rótulos do YOLO, com
normalização das coordenadas das caixas. Terceiro, a filtragem e o
balanceamento: serão descartadas imagens sem anotação e instâncias
degeneradas, e será avaliado o efeito do desbalanceamento entre classes,
dado que buracos são significativamente menos frequentes que trincas em
ambas as bases. Quarto, a divisão dos dados em treino, validação e teste
na proporção aproximada de 80\%, 10\% e 10\%, com a separação feita por
imagem de origem para evitar vazamento entre partições.

Serão aplicadas técnicas de aumento de dados sobre o conjunto de treino
(espelhamento horizontal e vertical, rotação, variação de escala, ajuste
de brilho e contraste e inserção de ruído), com o objetivo de expor o
modelo a diferentes condições de iluminação e ângulo de captura.
Seguindo a recomendação de Wanderley e Silva (2025), o aumento será
aplicado apenas ao treino, mantendo validação e teste inalterados, e uma
amostra das imagens transformadas será inspecionada visualmente para
garantir que as caixas delimitadoras permaneçam corretas após cada
transformação.

\subsection{Pré-processamento dos dados}
\label{sec:preprocessamento}

As etapas de preparação descritas na Seção~\ref{sec:dados} foram
materializadas em um único procedimento reprodutível, implementado no
\emph{script} \texttt{src/merge\_datasets.py}. O \emph{script} lê os dois
conjuntos de dados diretamente de seus arquivos de distribuição
compactados e produz, em uma única passagem, todos os artefatos
necessários ao treinamento: as imagens selecionadas, os rótulos no
formato YOLO, o arquivo de configuração \texttt{data.yml}, um manifesto
em CSV com uma linha por imagem exportada e um arquivo
\texttt{stats.json} com as contagens por fonte e por partição.

A opção por concentrar todo o pré-processamento em um único gerador, em
vez de encadear transformações sucessivas sobre um diretório de trabalho,
é deliberada. Ela garante que os cinco artefatos de saída sejam
mutuamente consistentes por construção, pois o manifesto nunca descreve
um conjunto que deixou de existir, e que todo o resultado seja
recuperável a partir de uma semente aleatória fixa, requisito de
reprodutibilidade já estabelecido na Seção~\ref{sec:treinamento}.
% OBS: na proposta original esta referencia apontava para a "Secao 4.5".
% Confira se o alvo pretendido era a secao de treinamento ou a de
% tecnologias (\ref{sec:tecnologias}).
Os conjuntos originais permanecem intactos, de modo que variantes do
conjunto de treinamento podem ser geradas e comparadas sem risco de perda
de dados.

As subseções seguintes descrevem as cinco decisões de pré-processamento
adotadas, na ordem em que são aplicadas pelo gerador.

\subsubsection{Restrição ao domínio de captura aérea}
\label{sec:restricao-aerea}

O Earendil destina-se à inspeção viária por drone, mas o conjunto
unificado construído a partir do RDD2022 e do UAV-PDD2023 não
representava essa perspectiva. Das 5.062 imagens reunidas na primeira
versão do conjunto, apenas 261, cerca de 5\%, haviam sido capturadas por
VANT; as demais provinham de câmeras embarcadas em automóveis e
motocicletas. As duas perspectivas não constituem variações menores de um
mesmo problema de visão, como sintetiza a Tabela~\ref{tab:perspectivas}.

\begin{table}[ht]
\centering
\caption{Diferenças entre as perspectivas de captura}
\label{tab:perspectivas}
\small
\setlength{\tabcolsep}{4pt}
\renewcommand{\arraystretch}{1.15}
\begin{tabular}{p{2.9cm} p{5.3cm} p{4.4cm}}
\hline
\textbf{Aspecto} & \textbf{Captura em nível de solo} &
\textbf{Captura por drone} \\
\hline
Atitude da câmera & Oblíqua, com o horizonte no quadro &
Nadir (perpendicular ao solo) \\
Cobertura da via & Terço inferior da imagem & Imagem inteira \\
Conteúdo da cena & Céu, edificações, tráfego, vegetação &
Apenas a textura do pavimento \\
Aparência do buraco & Pequeno, escorçado, distorcido pela perspectiva &
Quase planar, em escala consistente \\
Iluminação & Mista, com sombras projetadas intensas &
Comparativamente uniforme \\
\hline
\end{tabular}
\fonte{elaborado pelo autor, 2026.}
\end{table}

Um detector treinado majoritariamente sobre quadros oblíquos de câmera
veicular aprende uma geometria de projeção e um contexto de cena que não
ocorrem em imagens nadir de drone. Manter essa maioria no conjunto de
treinamento equivaleria a otimizar o modelo para uma condição de operação
que não é a do sistema proposto. O conjunto foi, por isso, restringido às
fontes aéreas.

O UAV-PDD2023 é integralmente capturado por drone e qualifica-se por
completo. O RDD2022, por sua vez, é organizado por país, e os países
diferem quanto à plataforma de captura: seis dos sete subconjuntos
(Índia, Japão, República Tcheca, Noruega, Estados Unidos e
\texttt{China\_MotorBike}) são veiculares, ao passo que
\texttt{China\_Drone} é aéreo, o que se confirma pela inspeção de suas
imagens: quadros de $512 \times 512$ pixels do pavimento vistos de cima,
sem horizonte e sem contexto de borda de via, visualmente compatíveis com
o material do UAV-PDD2023. A restrição não consiste, portanto, em
descartar o RDD2022, mas em preservar seu subconjunto aéreo e dispensar
os seis subconjuntos de nível de solo.

O filtro é aplicado no momento da indexação, e não na exportação. A
distinção tem consequência prática relevante: os seis arquivos excluídos
sequer chegam a ser abertos, apenas o da Noruega ocupa 10,6 GB, e a
geração completa do conjunto aéreo conclui-se em segundos. O
comportamento é controlado pela opção \texttt{-{}-drone-only},
desativada por padrão, de modo que o conjunto misto original permanece
reprodutível a partir do mesmo \emph{script}, condição necessária para
que a comparação entre os dois regimes de treinamento seja justa.

\subsubsection{Harmonização de taxonomias e conversão de formatos}
\label{sec:harmonizacao}

As duas bases nomeiam o buraco de maneiras distintas: o RDD2022 o anota
sob o código \texttt{D40} e o UAV-PDD2023 sob o rótulo
\texttt{Pothole}. Ambos são mapeados para a classe única
\texttt{0 Pothole} do conjunto unificado. As demais categorias de
patologia presentes nas bases, isto é, trincas longitudinais,
transversais, oblíquas e interligadas, além de remendos, são descartadas
em vez de reaproveitadas como classes auxiliares. A consequência dessa
decisão é que uma imagem que contenha apenas trincas não é removida do
conjunto: ela passa a constituir um exemplo negativo, isto é, uma amostra
de pavimento sem buraco, o que é coerente com o escopo delimitado na
Seção~\ref{sec:delimitacoes}, em que apenas a classe buraco é objeto de
avaliação.

As anotações originais, ambas no formato Pascal VOC, são convertidas para
o formato de rótulos do YOLO, com as coordenadas do centro, da largura e
da altura de cada caixa normalizadas em relação às dimensões da imagem.
Três cuidados foram incorporados à conversão:

\begin{itemize}
  \item \textbf{Dimensões obtidas da própria imagem.} A normalização
    utiliza a largura e a altura reais do arquivo JPEG decodificado, e
    não o campo \texttt{<size>} registrado no XML, que se mostrou
    inconsistente em parte das anotações do RDD2022. Normalizar por uma
    dimensão incorreta produziria caixas deslocadas de forma silenciosa,
    sem qualquer erro de execução.
  \item \textbf{Recorte às bordas e correção de orientação.} Coordenadas
    invertidas são reordenadas e caixas que extrapolam os limites da
    imagem são recortadas a eles.
  \item \textbf{Descarte de caixas degeneradas.} Caixas com menos de dois
    pixels em qualquer dos lados após o recorte são eliminadas, por não
    representarem objeto anotável e por degradarem o cálculo da função de
    perda.
\end{itemize}

Imagens para as quais não existe anotação correspondente no arquivo, bem
como arquivos truncados ou ilegíveis, são registrados em \emph{log} e
ignorados.

\subsubsection{Seleção de exemplos negativos}
\label{sec:negativos}

Todas as imagens que contêm ao menos um buraco são retidas. A elas
somam-se exemplos negativos na proporção de 0,33 negativo por positivo,
amostrados aleatoriamente sob a semente fixa. A presença controlada de
negativos expõe o detector a pavimento íntegro e a patologias que não são
buraco, reduzindo a tendência a predizer em excesso; a proporção é
mantida deliberadamente abaixo da paridade para não diluir o sinal
positivo em um conjunto já reduzido.

A razão é calculada dentro de cada fonte, e não sobre o agregado, de modo
que o equilíbrio entre positivos e negativos se verifique em cada base
isoladamente e não apenas na soma. Cada imagem negativa recebe um arquivo
de rótulo vazio, convenção pela qual o YOLO a interpreta como imagem sem
objetos, em vez de descartá-la por ausência de anotação.

\subsubsection{Partição em treino, validação e teste}
\label{sec:particao}

O UAV-PDD2023 distribui, junto das fotografias originais, cópias já
aumentadas de cada uma delas: espelhamentos horizontal e vertical e
rotação de $180^\circ$, identificados pelos prefixos \texttt{lr\_},
\texttt{tb\_} e \texttt{r180\_}, e recortes por quadrante, identificados
por sufixos como \texttt{\_top\_left}. Tratar esses arquivos como
amostras independentes seria um erro metodológico grave: cópias quase
idênticas de uma mesma fotografia se distribuiriam entre treino e
validação, e o modelo seria validado sobre imagens que já viu, inflando
artificialmente a métrica.

A partição é feita, por isso, sobre grupos e não sobre arquivos. Todas as
cópias derivadas de uma mesma fotografia de origem constituem um grupo e
são necessariamente atribuídas à mesma partição. Pela mesma razão, as
listas oficiais de partição distribuídas pelo UAV-PDD2023
(\texttt{ImageSets/Main}) não foram utilizadas: a verificação mostrou que
184 de seus 185 grupos de validação ocorrem também no conjunto de treino.
As imagens do RDD2022 são independentes entre si e cada uma constitui seu
próprio grupo.

A atribuição é ainda estratificada pelo par (fonte, presença de buraco),
de modo que ambas as fontes e ambas as condições estejam representadas em
validação e em teste, e emprega alocação por maiores restos para que
nenhuma partição seja silenciosamente reduzida a zero em estratos
pequenos. As proporções adotadas são de 80\% para treino, 10\% para
validação e 10\% para teste.

\subsubsection{Aumento de dados}
\label{sec:aumento}

O aumento de dados é aplicado exclusivamente ao conjunto de treino, em
tempo de execução, pelo \emph{pipeline} da biblioteca Ultralytics, que a
cada época sorteia transformações de espelhamento, escala, translação e
ajuste de matiz, saturação e brilho sobre as imagens carregadas,
propagando a transformação correspondente às caixas delimitadoras. Os
conjuntos de validação e teste não recebem qualquer transformação,
conforme a recomendação de Wanderley e Silva (2025) incorporada na
Seção~\ref{sec:dados}.

Optou-se por não gerar cópias aumentadas em disco durante o
pré-processamento. O aumento em tempo de execução expõe o modelo a uma
variação distinta a cada época, em vez de a um conjunto fixo de cópias,
e, o que é determinante neste caso, não cria novos arquivos cuja
filiação a uma fotografia de origem precisaria ser rastreada pelo
mecanismo de grupos descrito na Seção~\ref{sec:particao}, preservando a
integridade da partição.

\subsubsection{Conjunto resultante}
\label{sec:conjunto-resultante}

O conjunto aéreo resultante é gravado em \texttt{dataset/drone/},
preservado o conjunto misto em \texttt{dataset/merged/} para efeito de
comparação. Sua composição é apresentada nas
Tabelas~\ref{tab:composicao} e~\ref{tab:particoes}.

\begin{table}[ht]
\centering
\caption{Composição do conjunto aéreo, por fonte}
\label{tab:composicao}
\small
\setlength{\tabcolsep}{5pt}
\renewcommand{\arraystretch}{1.15}
\begin{tabular}{l r r r r l}
\hline
\textbf{Fonte} & \textbf{Imagens} & \textbf{Positivas} &
\textbf{Negativas} & \textbf{Caixas} & \textbf{Resolução} \\
\hline
UAV-PDD2023 & 176 & 132 & 44 & 195 & $2592 \times 1944$ \\
RDD2022 (\texttt{China\_Drone}) & 85 & 64 & 21 & 86 & $512 \times 512$ \\
\hline
\textbf{Total} & \textbf{261} & \textbf{196} & \textbf{65} &
\textbf{281} & \\
\hline
\end{tabular}
\fonte{elaborado pelo autor, 2026.}
\end{table}

\begin{table}[ht]
\centering
\caption{Distribuição por partição}
\label{tab:particoes}
\small
\setlength{\tabcolsep}{6pt}
\renewcommand{\arraystretch}{1.15}
\begin{tabular}{l r r r r}
\hline
\textbf{Partição} & \textbf{Imagens} & \textbf{Imagens positivas} &
\textbf{Caixas} & \textbf{Grupos} \\
\hline
Treino    & 210 & 159 & 228 & 120 \\
Validação &  28 &  21 &  28 &  14 \\
Teste     &  23 &  16 &  25 &  16 \\
\hline
\end{tabular}
\fonte{elaborado pelo autor, 2026.}
\end{table}

A restrição preserva a totalidade do material aéreo disponível: as 281
caixas de buraco correspondem a todas as que existem nas duas bases em
imagens capturadas por drone (195 no UAV-PDD2023 e 86 no
\texttt{China\_Drone}), de modo que nada se perdeu além do material de
nível de solo que o filtro se destinava a remover.

A integridade do conjunto gerado foi verificada por quatro checagens
automatizadas: que apenas fontes aéreas figuram no manifesto; que a
contagem total de caixas corresponde à esperada; que nenhum grupo de
aumento se distribui entre duas partições; e que todos os rótulos são
sintaticamente válidos, com identificador de classe igual a zero e
coordenadas normalizadas no intervalo adequado. Todas as checagens são
satisfeitas pelo conjunto em uso.

Como o diretório de dados não é versionado (os arquivos de distribuição e
os conjuntos gerados somam dezenas de gigabytes), o registro reprodutível
do conjunto é o próprio gerador, acompanhado da semente e dos parâmetros
registrados em \texttt{stats.json}.

\subsubsection{Limitações decorrentes do pré-processamento}
\label{sec:limitacoes-preproc}

Duas consequências da restrição ao domínio aéreo devem ser declaradas
junto de qualquer resultado obtido sobre este conjunto.

O conjunto é pequeno e sua diversidade efetiva é menor que a contagem de
imagens sugere. As 132 imagens positivas do UAV-PDD2023 derivam de apenas
31 fotografias distintas; as demais são espelhamentos, rotações e
recortes dessas 31. Somadas aos 64 quadros independentes do
\texttt{China\_Drone}, o conjunto contém aproximadamente 95 cenas de
buraco genuinamente independentes. Trata-se de um teto imposto pelas
bases de origem, e não pelo filtro: o UAV-PDD2023 contém 195 caixas de
buraco contra 5.398 de trinca transversal, e o \texttt{China\_Drone}
fornece 86, todas já incluídas.

Os conjuntos de validação e teste são reduzidos. A partição deixa 21
imagens positivas de validação, oriundas de 9 cenas independentes, e 16
imagens positivas de teste, oriundas de 10. Nessa escala, uma única
imagem desloca o mAP em vários pontos percentuais, de modo que a métrica
de validação oscilará apreciavelmente entre épocas e o critério de parada
antecipada repousará sobre sinal fraco. Os valores obtidos sobre o
conjunto de teste devem ser reportados com essa ressalva e não comparados
aos do conjunto misto como se houvessem sido medidos com precisão
equivalente.

Ambas as limitações são propriedades dos dados aéreos disponíveis
publicamente. Caso se mostrem impeditivas, três mitigações estão
previstas: elevar a proporção de negativos, aproveitando as cerca de
4.645 imagens aéreas sem buraco ainda não utilizadas nas duas bases;
adotar uma divisão mais larga, da ordem de 70/15/15; ou substituir a
partição única por validação cruzada sobre os grupos, o que produz
estimativa mais estável ao custo de múltiplos treinamentos.

\subsection{Treinamento e avaliação do modelo}
\label{sec:treinamento}

O detector será implementado com a arquitetura \textbf{YOLO26}, da
biblioteca Ultralytics, por três razões. A primeira é a inferência
nativamente ponta a ponta: o modelo produz as detecções finais
diretamente, sem a etapa de supressão de não-máximos exigida por versões
anteriores, o que reduz a latência e elimina uma fonte de divergência de
comportamento entre plataformas, vantagem operacional direta para um
serviço de inferência acionado pelo \emph{backend}. A segunda é a
orientação a dispositivos de borda e a execução eficiente em CPU,
característica que viabiliza a implantação do protótipo sem GPU dedicada
e habilita a evolução futura para inferência embarcada. A terceira são as
melhorias reportadas na detecção de objetos pequenos, decorrentes das
funções de perda ProgLoss e STAL, aspecto crítico em imagens aéreas, nas
quais o buraco ocupa fração reduzida do quadro (Ultralytics, 2026).

\begin{figure}[ht]
\centering
\includegraphics[width=\textwidth]{figuras/yolo26.png}
\caption{Comparação de desempenho entre o YOLO26 e arquiteturas
anteriores de detecção de objetos}
\label{fig:yolo26}
\fonte{Ultralytics (2026).}
\end{figure}

O treinamento empregará transferência de aprendizado a partir de pesos
pré-treinados em base de propósito geral, estratégia que reduz o tempo de
convergência e melhora a generalização em conjuntos de tamanho moderado.
Serão treinadas as variantes \emph{nano} e \emph{small} da arquitetura,
com resolução de entrada de 640 pixels a princípio, e os hiperparâmetros
de taxa de aprendizado, tamanho de lote e número de épocas serão
ajustados a partir do comportamento das curvas de perda no conjunto de
validação. Como linha de base comparativa, modelos YOLOv8n e YOLO11n
serão treinados sobre exatamente o mesmo conjunto e protocolo, de modo a
isolar o efeito da arquitetura. Os experimentos serão executados em
ambiente com aceleração por GPU, em Google Colab ou em máquina local
equivalente, com registro sistemático de configuração, sementes
aleatórias e artefatos de cada execução para assegurar
reprodutibilidade.

A avaliação do modelo empregará as métricas consagradas em detecção de
objetos (Padilla et al., 2020): mAP@50 e mAP@50-95, precisão, revocação e
F1, calculadas sobre o conjunto de teste, que não participa de nenhuma
etapa de treinamento. Serão medidos também o tempo de inferência por
imagem em CPU e em GPU e o tamanho do modelo em disco, indicadores
relevantes para a decisão de implantação. A análise dos erros será
qualitativa além de quantitativa, com inspeção das imagens que geraram
falsos positivos e falsos negativos, para identificar os confundidores
recorrentes já apontados na literatura, como remendos, manchas de óleo,
sombras intensas e bueiros.

\subsection{Tecnologias e justificativas}
\label{sec:tecnologias}

\textbf{Modelo e treinamento.} Python, com PyTorch e a biblioteca
Ultralytics, que fornece implementação de referência do YOLO26 e das
versões anteriores usadas como linha de base, além de rotinas
padronizadas de avaliação, condição para que a comparação entre
arquiteturas seja justa.

\textbf{Backend.} Python com FastAPI, escolhido pela validação automática
de dados de entrada, pelo suporte nativo a operações assíncronas (útil no
tratamento de envios concorrentes de imagens) e pela geração automática
de documentação da API, que facilita o desenvolvimento independente do
\emph{frontend}. Manter o \emph{backend} na mesma linguagem do modelo
elimina a necessidade de uma camada de comunicação entre processos em
linguagens distintas. A persistência usará MariaDB, por sua maturidade no
tratamento de tipos numéricos e temporais e pela disponibilidade de
extensões geoespaciais caso a evolução do projeto as demande, com acesso
via \emph{shell} (linha de comando).

\textbf{Frontend.} HTML, CSS e JavaScript, pela simplicidade de
implantação e pela ausência de etapa de compilação, suficientes para a
interface de consulta prevista.
% OBS: a versao original justificava o frontend pela "tipagem estatica",
% o que nao se aplica a HTML/CSS/JS puro. Ajuste conforme a stack real
% (se for usar TypeScript, vale retomar o argumento da tipagem).

\textbf{Infraestrutura e qualidade.} Todos os serviços serão empacotados
em contêineres Docker e orquestrados por Docker Compose, garantindo
reprodutibilidade do ambiente e simplicidade de implantação. Os testes
automatizados serão escritos com \emph{pytest} no \emph{backend}, e uma
esteira de integração contínua no GitHub Actions executará a suíte a cada
alteração.

\subsection{Testes e validação}
\label{sec:testes}

A verificação do sistema ocorrerá em quatro níveis. Os testes unitários
cobrirão as funções de normalização e renomeação de arquivos, validação
de metadados, cálculo de limiares e serialização das detecções. Os testes
de integração exercitarão o fluxo completo, da submissão da imagem à
consulta do registro pela API, com banco e repositório de arquivos
efêmeros, e verificarão os casos de erro: imagem corrompida, formato não
suportado, coordenadas ausentes ou inválidas e indisponibilidade do
serviço de inferência. Um teste de carga leve medirá o tempo de resposta
e a taxa de falhas sob envio concorrente de imagens, estabelecendo o
limite prático de operação do protótipo. Por fim, uma validação ponta a
ponta submeterá ao sistema um lote de imagens reais de drone não
pertencentes ao conjunto de treinamento, comparando as detecções
produzidas com a anotação manual das mesmas imagens.

\subsection{Critérios de sucesso}
\label{sec:criterios}

O projeto será considerado bem-sucedido se satisfizer os seguintes
critérios, verificáveis de forma objetiva:

\begin{itemize}
  \item \textbf{Funcional.} O fluxo completo opera de ponta a ponta, sem
    intervenção manual entre a submissão da imagem pela fonte externa e a
    exibição do registro no \emph{frontend}, com os quatro componentes
    integrados e implantáveis por um único comando.
  \item \textbf{Desempenho do modelo.} mAP@50 igual ou superior a 0,65,
    precisão igual ou superior a 0,75 e revocação igual ou superior a
    0,70 para a classe buraco no conjunto de teste, patamares situados
    acima dos relatados por Batista et al. (2025) e por Wanderley e Silva
    (2025) em contexto brasileiro, e abaixo dos obtidos em bases mais
    controladas, o que os torna simultaneamente ambiciosos e
    alcançáveis.
  \item \textbf{Desempenho do sistema.} Tempo mediano de processamento
    por imagem inferior a um segundo em CPU, medido entre o recebimento
    da requisição e a persistência do resultado.
  \item \textbf{Qualidade e usabilidade.} Cobertura de testes
    automatizados igual ou superior a 70\% nos módulos de \emph{backend}
    e pontuação SUS igual ou superior a 68, valor de referência para
    usabilidade adequada.
\end{itemize}

O não atingimento dos patamares de desempenho do modelo não invalida o
trabalho: nesse caso, a análise das causas (desbalanceamento de classes,
limitações das bases quanto a pavimentos urbanos brasileiros ou
inadequação da resolução de entrada ao tamanho aparente dos defeitos)
passa a constituir resultado da pesquisa, desde que os demais critérios
sejam satisfeitos.

% =====================================================================
\section{Cronograma}
\label{sec:cronograma}

As atividades estão distribuídas entre Projeto Final I, dedicado à
fundamentação, à modelagem, à preparação dos dados e aos primeiros
experimentos de treinamento, e Projeto Final II, voltado ao refinamento
do modelo, à construção do \emph{frontend}, à integração e à validação.
Nenhuma atividade havia sido executada até a elaboração desta proposta.


\newcommand{\x}{$\bullet$}

\begin{table}[ht]
\centering
\caption{Cronograma de atividades}
\label{tab:cronograma}
\footnotesize
\setlength{\tabcolsep}{3pt}
\renewcommand{\arraystretch}{1.2}
\begin{tabular}{p{5.8cm} *{10}{c}}
\hline
\textbf{Atividade} & \textbf{Ago} & \textbf{Set} & \textbf{Out} &
\textbf{Nov} & \textbf{Dez} & \textbf{Mar} & \textbf{Abr} & \textbf{Mai} &
\textbf{Jun} & \textbf{Jul} \\
\hline
Definição do problema                     & \x &    &    &    &    &    &    &    &    &    \\
Pesquisa de trabalhos relacionados        &    & \x &    &    &    &    &    &    &    &    \\
Pesquisa de datasets existentes           &    & \x &    &    &    &    &    &    &    &    \\
Projeto da Arquitetura inicial do Sistema &    &    & \x & \x &    &    &    &    &    &    \\
Elaboração do texto da proposta de TCC    &    &    &    & \x & \x &    &    &    &    &    \\
Treinamento do modelo                     &    &    &    &    &    & \x &    &    &    &    \\
Implementação do \emph{backend}           &    &    &    &    &    &    & \x &    &    &    \\
Implementação do \emph{frontend}          &    &    &    &    &    &    & \x &    &    &    \\
Integração dos módulos do sistema         &    &    &    &    &    &    &    & \x &    &    \\
Realizar testes ponta a ponta no sistema &    &    &    &    &    &    &    & \x & \x  &    \\
Redação da monografia                     &    &    &    &    & \x & \x & \x & \x & \x & \x \\
\hline
\end{tabular}
\fonte{elaborado pelo autor, 2026.}
\end{table}

% =====================================================================
\section*{Referências}
\addcontentsline{toc}{section}{Referências}

\begin{referencias}

\item Alzamzami, O., Babour, A., Baalawi, W., \& Al Khuzayem, L. (2024).
PDS-UAV: A deep learning-based pothole detection system using unmanned
aerial vehicle images. \emph{Sustainability}, 16(21), 9168.
\url{https://doi.org/10.3390/su16219168}

\item Arya, D., Maeda, H., Sekimoto, Y., Omata, H., Ghosh, S. K.,
Toshniwal, D., et al. (2024). RDD2022: A multi-national image dataset for
automatic road damage detection. \emph{Geoscience Data Journal}.
\url{https://doi.org/10.1002/gdj3.260}

\item Batista, N. S., Oliveira, V. C., Souza, T. I. A., Castro, M. T.,
Gomes, D. G., Bertoncini, B. V., \& Branco, V. T. F. C. (2025). Detecção
automática de defeitos em pavimentos: uma abordagem com \emph{deep
learning} via YOLOv8n. Em \emph{Anais do XIX Brazilian e-Science
Workshop} (pp. 57-64). Porto Alegre: SBC.
\url{https://doi.org/10.5753/bresci.2025.248214}

\item Carvalho, R. B., Guedes, E. B., \& Figueiredo, C. M. S. (2022).
Detecção inteligente de falhas em pavimentações asfálticas com redes
neurais convolucionais regionais. Em \emph{Anais do III Workshop
Brasileiro de Cidades Inteligentes} (pp. 119-130). Porto Alegre: SBC.
\url{https://doi.org/10.5753/wbci.2022.222795}

\item CNT (2025). \emph{Pesquisa CNT de Rodovias 2025}. Brasília:
Confederação Nacional do Transporte.

\item DNIT (2003). \emph{DNIT 005/2003-TER: Defeitos nos pavimentos
flexíveis e semirrígidos. Terminologia}. Rio de Janeiro: Departamento
Nacional de Infraestrutura de Transportes.

\item DNIT (2003). \emph{DNIT 008/2003-PRO: Levantamento visual contínuo
para avaliação da superfície de pavimentos flexíveis e semirrígidos.
Procedimento}. Rio de Janeiro: Departamento Nacional de Infraestrutura de
Transportes.

\item Guan, T., Cai, J., Wang, Y., Yang, W., Chang, X., \& Han, Y.
(2025). Pavement pothole detection system based on deep learning and
binocular vision. \emph{Journal of Traffic and Transportation Engineering
(English Edition)}, 12(4), 1100-1123.
\url{https://doi.org/10.1016/j.jtte.2024.08.001}

\item Kassem, A. R. K., Matrakas, M. D., \& Bussador, A. (2024).
Algoritmo de detecção de rachaduras no concreto asfáltico betuminoso por
meio de visão computacional. Em \emph{Anais do XXI Congresso
Latino-Americano de Software Livre e Tecnologias Abertas} (pp. 356-359).
Porto Alegre: SBC.
\url{https://doi.org/10.5753/latinoware.2024.245330}

\item Monte, D. M. F. M. do, Goldschmidt, R. R., \& Santos, M. dos
(2025). Identificação e classificação de patologias em pavimentos
rodoviários flexíveis e rígidos utilizando YOLOv8. Em \emph{Proceedings
of the 31st Brazilian Symposium on Multimedia and the Web} (pp.
293-301). Porto Alegre: SBC.
\url{https://doi.org/10.5753/webmedia.2025.16131}

\item Padilla, R., Netto, S. L., \& da Silva, E. A. B. (2020). A survey
on performance metrics for object-detection algorithms. Em
\emph{Proceedings of the International Conference on Systems, Signals and
Image Processing (IWSSIP)} (pp. 237-242). Niterói: IEEE.

\item Redmon, J., Divvala, S., Girshick, R., \& Farhadi, A. (2016). You
only look once: Unified, real-time object detection. Em
\emph{Proceedings of the IEEE Conference on Computer Vision and Pattern
Recognition (CVPR)} (pp. 779-788). Las Vegas: IEEE.

\item Ultralytics (2026). \emph{Ultralytics YOLO26}. Recuperado de
\url{https://docs.ultralytics.com/models/yolo26/}

\item Wanderley, B. P., \& Silva, P. B. da (2025). Análise de patologia
de pavimentos através de fotos capturadas por drone utilizando IA. Em
\emph{Anais da XXI Jornada de Iniciação Científica}. São Paulo:
Universidade Presbiteriana Mackenzie.

\item Yan, H., \& Zhang, J. (2023). UAV-PDD2023: A benchmark dataset for
pavement distress detection based on UAV images. \emph{Data in Brief},
51, 109692. \url{https://doi.org/10.1016/j.dib.2023.109692}

\end{referencias}

\end{document}

